\documentclass[11pt,a4paper]{article}
\usepackage[margin=1in]{geometry}
\usepackage{hyperref}
\usepackage{amsmath}
\usepackage{sectsty}
\usepackage{etoolbox}

% Make section titles bold and slightly larger
\sectionfont{\large\bfseries}

\title{Learning by Osmosis: Acquiring Tacit Knowledge Through Immersive Participation in Communities of Practice and Information-Rich Environments}
\author{groda \\ Independent Researcher}
%\date{\today}
\date{April 4, 2026}

\begin{document}

\maketitle

\begin{abstract}
In an era of information abundance, the most valuable forms of knowledge---tacit, contextual, and difficult to codify---often remain inaccessible to outsiders who rely solely on deliberate search. ``Learning by osmosis,'' alternatively framed as ``learning by focusing,'' refers to the process whereby sustained immersion in a specific domain (online communities, physical laboratories, research institutes, innovation labs, departmental meetings, working groups, or conversational interfaces with AI) creates conditions for the passive-yet-active absorption of tacit knowledge. Rather than actively hunting for information, the learner curates an environment of focused attention; relevant insights, tools, norms, and intuitions then flow toward them through repeated exposure, casual interactions, and micro-serendipities. Drawing on theories from situated learning, knowledge management, and information retrieval, this article examines the phenomenon's theoretical foundations, the critical distinction between explicit and tacit knowledge, underlying mechanisms, empirical factors influencing its effectiveness (``osmosis rate''), expanded real-world examples across physical and hybrid settings, and practical implications. It argues that learning by osmosis is a deliberate, highly efficient strategy for mastering complex domains in both traditional and digital contexts.
\end{abstract}

\section{Defining the Concept}
\textbf{Learning by osmosis} describes the absorption of knowledge that occurs when one ``hangs around'' inside a field: listening to informal conversations, scrolling high-signal feeds, participating in routine activities, or engaging in low-stakes experimentation. Information arrives unbidden because the learner has positioned themselves within the stream. The complementary term \textbf{``learning by focusing''} highlights the intentional narrowing of attention to one domain---whether a Discord server, a research lab, a departmental meeting series, or a working group---thereby engineering a high-signal environment.

This approach contrasts sharply with explicit, search-driven learning. Search engines, databases, and curated lists deliver \emph{explicit} knowledge (facts, tutorials, papers). Osmotic or focused learning primarily yields \emph{tacit} knowledge: the unspoken practices, intuitions, shortcuts, and ``vibe'' that outsiders rarely discover (Polanyi, 1966). In information retrieval (IR) terms, it aligns with \emph{information encountering}---the accidental discovery of useful information during the pursuit of something else (Erdelez, 1997).

\section{Explicit versus Tacit Knowledge: A Core Distinction}
The effectiveness of learning by osmosis rests on a fundamental distinction, first articulated by philosopher Michael Polanyi, between \emph{explicit} and \emph{tacit} knowledge (Polanyi, 1966). Explicit knowledge is codified, declarative, and easily transmitted through language, documents, or digital media. It consists of facts, rules, procedures, formulas, and step-by-step instructions that can be searched, stored, and retrieved at will---think of a tutorial on how to fine-tune a language model, a published research paper, or a database entry listing available AI tools. Because it is articulated and decontextualized, explicit knowledge lends itself perfectly to traditional search-driven or classroom-style learning. Anyone with the right query can, in principle, access it regardless of prior immersion in the field.

Tacit knowledge, by contrast, is personal, context-dependent, and notoriously difficult to articulate fully. It encompasses the ``know-how,'' intuitions, subtle judgments, and embodied skills that experts possess but cannot easily put into words: sensing when a model is about to diverge, recognizing the unspoken cultural norms of a research lab, or instinctively choosing the right workflow for ``vibe coding'' based on years of pattern recognition. Polanyi's famous dictum---``we can know more than we can tell''---captures this gap. Tacit knowledge is acquired not through reading or direct instruction but through prolonged, situated experience: observing, imitating, experimenting, and participating in the social practices of a community. It cannot be fully transferred via manuals or search results; it must be internalized through repeated exposure and social interaction.

This distinction explains why outsiders often feel overwhelmed or lost when searching for domain-specific resources, while insiders effortlessly discover high-value insights. Explicit knowledge is abundant and noisy; tacit knowledge is scarce and environment-bound. Learning by osmosis (or focusing) is therefore the primary mechanism for acquiring the latter. By immersing themselves in a rich setting, learners do not merely accumulate facts---they internalize the unarticulated expertise that turns novices into competent practitioners. Subsequent theoretical frameworks build directly on this foundation, showing \emph{how} tacit knowledge is socialized and encountered in practice.

\section{Theoretical Foundations}
Building on the explicit--tacit distinction, learning by osmosis (or learning by focusing) rests on four interlocking frameworks from philosophy, the sociology of learning, knowledge management, and information retrieval.

\begin{itemize}
\item \textbf{Tacit knowledge} (Polanyi, 1966): As elaborated above, much of what experts know cannot be fully articulated and is acquired through prolonged immersion and participation rather than explicit instruction.
\item \textbf{Communities of Practice (CoP) and situated learning} (Lave \& Wenger, 1991): Knowledge resides within social groups sharing a common enterprise. Newcomers enter via \textbf{legitimate peripheral participation}---lurking, observing, and gradually contributing---internalizing the community's tacit norms and practices.
\item \textbf{SECI model of knowledge creation} (Nonaka \& Takeuchi, 1995): The \emph{socialization} mode (tacit-to-tacit transfer through shared experience) is the core of osmotic learning. It occurs in joint activities, informal interactions, and physical or virtual proximity---precisely the ``hanging around'' described here.
\item \textbf{Information encountering and serendipitous retrieval} (Erdelez, 1997; Toms, 2000): From the IR literature, these frameworks explain how immersion in information-rich environments generates constant mini-serendipities. Focused attention increases the probability of noticing, stopping at, and exploiting unexpected but relevant information.
\end{itemize}

Together, these theories show that immersion is not passive luck but the deliberate creation of conditions for ongoing, micro-level serendipitous encounters.

\section{Mechanisms at Work}
The process of learning by osmosis operates through several interconnected mechanisms that transform sustained immersion into reliable knowledge acquisition. These mechanisms bridge the gap between explicit search and tacit internalization, explaining how focused participation in a community or environment reliably generates the flow of otherwise invisible insights.

\begin{enumerate}
\item \textbf{Environmental curation (``focusing'')} : Restricting attention to one domain reduces noise and raises signal density.
\item \textbf{Social and experiential osmosis} : Casual conversations, shared artifacts (screenshots, demos, meeting notes), and repeated exposure transmit tacit norms faster than formal documents.
\item \textbf{Repeated low-stakes exposure} : Trying small experiments or attending routine meetings builds pattern recognition and intuition.
\item \textbf{Mini-serendipities via information encountering} : The longer one remains in the stream, the higher the likelihood that precisely the needed insight appears at the right moment (Erdelez, 1997; Toms, 2000).
\end{enumerate}

Physical or hybrid environments (labs, departments, working groups) add embodied cues and spontaneous ``corridor teaching,'' while digital ones scale access through persistent feeds and archives. Structured elements such as recurring meeting series, internal wikis, or working groups act as ritualized channels for socialization in the SECI sense.

\section{Empirical Insights: Factors Influencing the ``Osmosis Rate''}
Empirical research has identified and quantified several key factors that determine the effectiveness of tacit knowledge transfer within communities of practice and other information-rich environments---what we might informally refer to as the ``osmosis rate.'' These factors have been rigorously evaluated through quantitative methods such as structural equation modeling (SEM) and qualitative analyses of CoP dynamics, confirming their predictive power across research groups, higher-education institutions, and professional settings.

A prominent study by Alves et al.\ (2022) applied partial least squares structural equation modeling (PLS-SEM) to survey data from 255 professors and researchers in a Brazilian public higher-education institution. The model accounted for 27.3\% of the variance in tacit knowledge sharing ($R^2 = 0.273$) and identified three significant positive predictors: (1) \textbf{individual factors} ($\beta = 0.320$, $p < 0.001$; strongest predictor, encompassing awareness, critical thinking, time management, shared language, and willingness to share); (2) \textbf{organizational structure factors} ($\beta = 0.211$, $p < 0.01$; including proximity, relational networks, communication systems, and supportive infrastructure); and (3) \textbf{knowledge management strategy factors} ($\beta = 0.147$, $p < 0.05$; such as mentoring programs, socialization routines, and explicit encouragement of CoPs). Organizational culture alone showed no significant effect ($\beta = 0.004$, $p > 0.05$). These results underscore that structural supports and targeted strategies outweigh vague cultural openness in facilitating osmotic processes.

Complementary qualitative and mixed-methods research on CoPs further illuminates the mechanisms that elevate the osmosis rate. Pyrko et al.\ (2017) highlight ``thinking together''---mutual guidance through storytelling, joint problem-solving, and reflective dialogue---as a central driver that converts legitimate peripheral participation into deep tacit internalization. Broader syntheses of CoP evidence (Wenger-Trayner \& Wenger-Trayner, 2024; James-McAlpine et al., 2023) consistently identify success factors such as fluid boundaries, self-selection of passionate members, a regular rhythm of interaction, distributed leadership, and avoidance of excessive hierarchy. Barriers repeatedly documented include insufficient dedicated time, geographic or temporal distance, and over-formalization that suppresses spontaneous encounters (Alves et al., 2022).

Collectively, these empirically grounded findings demonstrate that the osmosis rate is neither random nor purely cultural; it is systematically shaped by identifiable, actionable levers at the individual, structural, and strategic levels. Environments that deliberately combine sustained proximity (or persistent digital presence), psychological safety, and low-stakes interaction channels consistently achieve higher rates of tacit knowledge absorption.

\section{Concrete Examples Across Physical, Organizational, and Hybrid Settings}
\textbf{Online communities in artificial intelligence and machine learning}  
In the rapidly evolving domain of artificial intelligence and machine learning---characterized by swift technological progress, intense commercial competition, and a consequent proliferation of promotional material---participants immersed in specialized online networks routinely acquire tacit insights unavailable through conventional search. Insiders active in X/Twitter circles, Discord servers (e.g., Latent Space, EleutherAI), and dedicated subreddits discover niche tools and workflows---such as custom ComfyUI pipelines, specialized fine-tuning services for scientific datasets, and ``vibe coding'' assistants---through casual mentions, shared demonstrations, and threaded discussions long before these resources surface in curated lists or formal publications. Attempts by outsiders to locate comparable resources via search engines or review articles typically encounter an overwhelming volume of marketing-driven content, outdated recommendations, and commercial noise, illustrating the inherent limitations of explicit knowledge channels in dynamic, economically saturated fields.

\textbf{Physical and institutional settings}: Medical residents absorb diagnostic intuition during ward rounds and corridor conversations far more effectively than from textbooks (Carsten Conner, 2018). PhD students in research institutes gain meta-knowledge about journals, reviewers, and experimental tricks through lab meetings, coffee breaks, and journal clubs. Innovation labs and working groups in companies or research institutes facilitate rapid tacit transfer via daily stand-ups and project reviews. Structured organizational routines---such as recurring meeting series, internal wikis, or working groups---function as socialization channels, surfacing insights that broad searches would miss.

\textbf{Scientific communities and departments}: Field science apprenticeships (e.g., immersive programs placing novices alongside experts) demonstrate how authentic, embodied participation in communities of practice accelerates tacit skill acquisition (Carsten Conner, 2018). Research institutes provide the classic ``participant observation'' environment: long-term hanging around turns peripheral observers into full contributors.

\textbf{Emerging hybrid: Chatting casually with AI}: Conversational large language models (e.g., daily interactions with systems like Grok or Claude) create a personalized, always-available ``community of one.'' Casual prompting about a workflow, a paper, or a problem often surfaces unexpected connections, analogies, or practical shortcuts---functioning as a form of simulated socialization. The back-and-forth dialogue mimics legitimate peripheral participation: the user ``focuses'' on the topic, and the AI supplies contextually relevant tacit-like insights on demand. This represents a novel extension of osmotic learning, blending human intentionality with machine-scale information encountering.

\section{Benefits, Limitations, and Practical Guidance}
While the theoretical, empirical, and mechanistic foundations of learning by osmosis underscore its power for acquiring hard-to-codify tacit knowledge, realizing these gains in practice requires both individual discipline and organizational stewardship. Osmotic learning is not automatic; its effectiveness depends on deliberate choices about where and how one immerses oneself, as well as on the design of the environments themselves. Evidence-based practical guidance, drawn from communities-of-practice research (Lave \& Wenger, 1991; Wenger-Trayner \& Wenger-Trayner, 2024), therefore addresses two complementary perspectives: strategies for individual learners and strategies for those who design and sustain osmotic environments.

\textbf{Benefits} include accelerated tacit expertise, stronger intuition, robust professional networks, and an almost effortless sense of mastery that search-driven approaches rarely achieve. Insiders gain rapid pattern recognition, contextual judgment, and access to the ``vibe'' of a field---advantages that translate into faster innovation, better problem-solving, and deeper professional identity (Pyrko et al., 2017; Wenger-Trayner \& Wenger-Trayner, 2024).

\textbf{Limitations} involve potential echo chambers, slower initial onboarding, and risk of missing cross-disciplinary insights. Remote or hybrid work can disrupt traditional osmotic channels unless deliberately compensated (Gratton, 2021). Over-formalization or excessive hierarchy can also stifle the spontaneous encounters that drive tacit transfer (Alves et al., 2022).

\subsection{Practical guidance for individual learners}
Identification of suitable environments constitutes a foundational step. Effective settings are characterized by high signal density, including regular interaction rhythms (weekly meetings, active Discord channels, recurring lab seminars), passionate self-selected members, fluid boundaries that support legitimate peripheral participation, and visible exchanges of value (shared demonstrations, storytelling, or joint problem-solving). Learners may assess potential communities through short trial periods of observation or review of public archives and recent interactions before committing. Sustained engagement---typically daily or weekly peripheral participation over a minimum of three to six months---is essential: this involves systematic reading of discussions, attendance at routine gatherings, and weekly experimentation with at least one suggested practice. Maintenance of a simple serendipity log documenting unexpected discoveries further reinforces habit formation and enables reflection on progress.

To mitigate risks of tunnel vision or echo chambers, individuals should incorporate deliberate boundary-spanning activities. Wenger-Trayner and Wenger-Trayner (2024) advocate periodic engagement with adjacent or contrasting communities (e.g., maintaining one primary domain focus alongside lighter exposure to a cross-disciplinary feed or secondary working group). Scheduling monthly exploration periods for external resources or events outside the core domain helps preserve the depth of focused immersion while introducing fresh perspectives and reducing groupthink. Environments lacking psychological safety---where naive questions or tentative ideas are discouraged---should be abandoned in favor of more conducive settings.

\subsection{Practical guidance for managers, research institute heads, meetup organizers, and community builders}
Leaders play a pivotal role in engineering environments that maximize the osmosis rate. Key design principles include provision of dedicated time and physical or virtual proximity, establishment of low-stakes interaction channels (recurring meeting series, internal wikis, informal forums), and cultivation of psychological safety to enable storytelling, joint problem-solving, and reflective dialogue (Alves et al., 2022; Pyrko et al., 2017). Designation of ``community gardeners''---facilitators who sustain interaction rhythms, welcome newcomers, and preserve fluid boundaries without imposing rigid hierarchy---further supports organic growth. A balanced mix of stabilizing routines (regular stand-ups) and stimulating activities (guest demonstrations, cross-team challenges) maintains engagement, while visible senior-leadership sponsorship (resource allocation, public recognition of contributions, and protection of participation time) signals legitimacy and sustains long-term momentum (Wenger-Trayner \& Wenger-Trayner, 2024).

Attraction and retention of participants are best achieved by beginning with a small core of intrinsically motivated insiders who share a genuine domain interest, then expanding organically through demonstration of early value (success stories, shared artifacts, and observable impact). Lowering barriers to entry---transparent onboarding processes, multiple levels of participation (from lurking to co-creation), and asynchronous options---enhances accessibility. Clear articulation of the community's domain focus and emphasis on reciprocal learning attract self-selected members. In research institutes or companies, integration of osmotic channels into existing workflows (e.g., embedding working-group time within project calendars) rather than treating them as optional extras promotes sustained participation. The outcome is a self-reinforcing ecosystem in which tacit knowledge flows reliably, transforming peripheral participants into full contributors and elevating collective expertise.

\section{Conclusion}
Learning by osmosis---reframed as learning by focusing---is an ancient yet powerfully relevant strategy for acquiring tacit knowledge in an information-saturated world. By drawing on communities of practice (Lave \& Wenger, 1991), the SECI socialization mode (Nonaka \& Takeuchi, 1995), and information encountering (Erdelez, 1997; Toms, 2000), we see that the highest-leverage move is often not broader searching but deeper immersion. Whether in physical research institutes, departmental meetings, online circles, or casual AI conversations, the key is to stand in the right place long enough for the field itself to teach us. In doing so, we transform serendipity from rare accident into reliable, everyday occurrence.

\section*{References}
\begin{thebibliography}{11}

\bibitem{alves2022}
Alves, R. B. C., \& Pinheiro, P. (2022). Factors influencing tacit knowledge sharing in research groups in higher education institutions. \emph{Administrative Sciences}, 12(3), 89. \href{https://doi.org/10.3390/admsci12030089}{https://doi.org/10.3390/admsci12030089}

\bibitem{carstenconner2018}
Carsten Conner, L. D. (2018). Tacit knowledge and girls' notions about a field science community of practice. \emph{International Journal of Science Education, Part B}, 8(2), 100--118.

\bibitem{erdelez1997}
Erdelez, S. (1997). Information encountering: A conceptual framework for accidental information discovery. In \emph{Proceedings of the 1997 ASIS Annual Meeting}.

\bibitem{gratton2021}
Gratton, L. (2021). The problem with losing ``osmosis learning''. \emph{BBC Worklife}, October 4.

\bibitem{jamesmcalpine2023}
James-McAlpine, J., et al. (2023). Exploring the evidence base for communities of practice in health research and translation. \emph{Health Research Policy and Systems}.

\bibitem{lavewenger1991}
Lave, J., \& Wenger, E. (1991). \emph{Situated learning: Legitimate peripheral participation}. Cambridge University Press.

\bibitem{nonakatakeuchi1995}
Nonaka, I., \& Takeuchi, H. (1995). \emph{The knowledge-creating company: How Japanese companies create the dynamics of innovation}. Oxford University Press.

\bibitem{polanyi1966}
Polanyi, M. (1966). \emph{The tacit dimension}. University of Chicago Press.

\bibitem{pyrko2017}
Pyrko, I., Dörfler, V., \& Eden, C. (2017). Thinking together: What makes communities of practice work? \emph{Research Policy}, 46(5), 938--954.

\bibitem{toms2000}
Toms, E. G. (2000). Serendipitous information retrieval. In \emph{Proceedings of the First DELOS Network of Excellence Workshop on Information Seeking, Searching and Querying in Digital Libraries}.

\bibitem{wengertayner2024}
Wenger-Trayner, E., \& Wenger-Trayner, B. (2024). \emph{Communities of practice guidebook} (2nd ed.).

\end{thebibliography}

\end{document}
